\subsection{Character of a Representation}

Let E be a left A-module of finite dimension over K. The character of E or trace of E, denoted by \(Tr_{E}\) , is the linear form \(a \mapsto \mathrm{Tr}(a_{\mathrm{E}})\) . Let \((e_{1}, \ldots, e_{n})\) be a basis of E and \((e_{1}^{*}, \ldots, e_{n}^{*})\) its dual basis. By definition, we have the relation

\[
\mathrm{Tr} _ {\mathrm{E}} = \sum_ {i = 1} ^ {n} c _ {\mathrm{E}} (e _ {i}, e _ {i} ^ {*}).\tag{7}
\]

The linear form \(Tr_{E}\) belongs to \(\Theta_{\mathrm{E}}(\mathrm{A})\) . We have \(Tr_{E} = Tr_{E'}\) if E and \(E'\) are two isomorphic A-modules. So \(Tr_{E}\) depends only on the isomorphism class of E.

Let \(\pi\) be a linear representation of A in a vector space V of finite dimension over K. The trace of \(\pi\) , or sometimes character of \(\pi\) , is defined as the character of the A-module \(V_{\pi}\) . We denote it by \(\mathrm{Tr}_{\pi}\) . If \((e_1, \ldots, e_n)\) is a basis of V over K and if \((\pi_{ij}(a))\) is the matrix of \(\pi(a)\) with respect to this basis, then we have

\[
\operatorname{Tr} _ {\pi} (a) = \sum_ {i = 1} ^ {n} \pi_ {i i} (a).\tag{8}
\]

The linear form \(\mathrm{Tr}_{\pi}\) belongs to \(\Theta_{\pi}(\mathrm{A})\) .

PROPOSITION 5. --- Let S be a simple A-module of finite dimension over K, let D be the opposite field of the commutant of S, and let Z be the center of D. The following properties are equivalent:

\begin{enumerate}
\def\labelenumi{(\roman{enumi})}
\item
  The character \(Tr_{S}\) of S is not zero.
\item
  There exists an element \(d \in \mathrm{D}\) such that \(\operatorname{Tr}_{\mathrm{D} / \mathrm{K}}(d) \neq 0\) .
\item
  The extension \(\mathbf{Z}\) of \(\mathbf{K}\) is separable, and the characteristic \(p\) of \(\mathbf{K}\) does not divide the degree \([\mathrm{D}:\mathrm{Z}]\) .
\end{enumerate}

The right D-vector space S is finite-dimensional. Let \((e_{1},\ldots,e_{n})\) be a basis of S over D, and let u be an element of \(\operatorname{End}_{\mathrm{D}}(\mathrm{S})\) . Let \((d_{ij})\) be the matrix of u with respect to the basis \((e_{1},\ldots,e_{n})\) . We have \(u(e_{j})=\sum_{i=1}^{n}e_{i}d_{ij}\) for \(j\in[1,n]\) . Denote by \(u_{K}\) the mapping u viewed as an endomorphism of the K-vector space S and by \((u_{ij})\) the matrix of \(u_{K}\) with respect to the decomposition \((e_{i}\mathrm{D})\) of the K-vector space S as a direct sum (II, §10, No.~5, p.~346). We have \(\operatorname{Tr}(u_{\mathrm{K}})=\sum_{i}\operatorname{Tr}(u_{ii})\) . Moreover, \(u_{ii}\) is the endomorphism of the K-vector space \(e_{i}D\) defined by \(u_{ii}(e_{i}d)=e_{i}d_{ii}d\) for \(d\in D\) , so its trace is equal to \(\operatorname{Tr}_{\mathrm{D}/\mathrm{K}}(d_{ii})\) . We have thus proved the equality

\[
\mathrm{Tr} (u _ {\mathrm{K}}) = \mathrm{Tr} _ {\mathrm{D/K}} \bigl (\sum_ {i} d _ {i i} \bigr).\tag{9}
\]

By Burnside's theorem (VIII, p.~83, Corollary 1 of Proposition 4), the mapping \(a \mapsto a_{\mathrm{S}}\) from A to \(\mathrm{End}_{\mathrm{D}}(\mathrm{S})\) is surjective. The equivalence of properties (i) and (ii) therefore follows from formula (9).

If the extension Z of K is separable, then there exists an element \(d \in Z\) such that \(\operatorname{Tr}_{\mathrm{Z/K}}(d) \neq 0\) (V, §8, No.~2, p.~49, Corollary of Proposition 1). Moreover, we have \(\operatorname{Tr}_{\mathrm{D/K}}(d) = [\mathrm{D} : \mathrm{Z}] \operatorname{Tr}_{\mathrm{Z/K}}(d)\) (III, §9, No.~4, p.~548, Corollary); consequently, if p does not divide {[}D : Z{]}, then we have \(\operatorname{Tr}_{\mathrm{D/K}}(d) \neq 0\) . This proves that (iii) implies (ii).

If the extension \(\mathrm{Z}\) of \(\mathrm{K}\) is not separable, then we have \(\operatorname{Tr}_{\mathrm{Z / K}}(x) = 0\) for every \(x \in \mathrm{Z}\) , and therefore \(\operatorname{Tr}_{\mathrm{D / K}}(d) = \operatorname{Tr}_{\mathrm{Z / K}}(\operatorname{Tr}_{\mathrm{D / Z}}(d)) = 0\) for every \(d \in \mathrm{D}\) .

Now suppose that p divides the degree \([D:Z]\) . It then divides the reduced degree of D over Z. For every \(d \in D\) , we have \(\mathrm{Tr}_{\mathrm{D}/\mathrm{Z}}(d) = 0\) (VIII, p.~344, Remark), and therefore \(\mathrm{Tr}_{\mathrm{D}/\mathrm{K}}(d) = \mathrm{Tr}_{\mathrm{Z}/\mathrm{K}}(\mathrm{Tr}_{\mathrm{D}/\mathrm{Z}}(d)) = 0\) . This completes the proof of the implication (ii) \(\Rightarrow\) (iii).

COROLLARY. --- If the field \(\mathbf{K}\) is perfect, then properties (i) through (iii) hold.

Since the field is perfect, the extension Z of K is separable (V, §15, No.~5, p.~125, Theorem 3). Property (iii) then follows from Corollary 3 of VIII, p.~323.

PROPOSITION 6. --- Let \(\mathcal{S}_0\) be the set of classes of simple A-modules of finite dimension over K with nonzero trace. The family of linear forms \((\mathrm{Tr}_{\mathrm{S}})_{\mathrm{S}\in \mathcal{S}_0}\) is free over K.

Let \(\mathrm{F}\) be a finite subset of \(\mathcal{S}_0\) , and let \(\mathrm{S}\) be an element of \(\mathrm{F}\) . By assumption, there exists an element \(a \in \mathrm{A}\) such that \(\operatorname{Tr}_{\mathrm{S}}(a) \neq 0\) . By Corollary 1 of Proposition 4 (VIII, p.~83), there exists an element \(b \in \mathrm{A}\) such that \(b_{\mathrm{S}} = a_{\mathrm{S}}\) and \(b_{\mathrm{T}} = 0\) for every \(\mathrm{T} \in \mathrm{F} - \{\mathrm{S}\}\) . We have \(\operatorname{Tr}_{\mathrm{S}}(b) \neq 0\) and \(\operatorname{Tr}_{\mathrm{T}}(b) = 0\) for \(\mathrm{T} \in \mathrm{F} - \{\mathrm{S}\}\) . The family \((\operatorname{Tr}_{\mathrm{S}})_{\mathrm{S} \in \mathrm{F}}\) is therefore free. Proposition 6 follows.

Remark. --- Let \(S_{K}\) be the set of classes of simple A-modules of finite dimension over K. Proposition 6 also follows from the fact that the sum of the \(\Theta_{\mathrm{S}}(\mathrm{A})\) , for S running through \(S_{K}\) , is direct.

Let

\[
0 \longrightarrow \mathrm{E} ^ {\prime} \longrightarrow \mathrm{E} \longrightarrow \mathrm{E} ^ {\prime \prime} \longrightarrow 0
\]

be an exact sequence of A-modules, all finite-dimensional over K. By Proposition 1 of III, §9, No.~2, p.~543, we have \(\mathrm{Tr}_{\mathrm{E}} = \mathrm{Tr}_{\mathrm{E}'} + \mathrm{Tr}_{\mathrm{E}''}\) . By the definition of the Grothendieck group \(\mathrm{R}_{\mathrm{K}}(\mathrm{A})\) (VIII, p.~194) and its universal property (VIII, p.~186, Proposition 4), there exists a homomorphism of additive groups \(\theta\) from \(\mathrm{R}_{\mathrm{K}}(\mathrm{A})\) to \(\mathrm{A}^*\) characterized by the relation \(\theta([\mathrm{E}]) = \mathrm{Tr}_{\mathrm{E}}\) for every A-module E of finite dimension over K. In particular, the trace of a semisimplification of E is equal to that of E. We deduce from \(\theta\) a K-linear mapping \(\theta_{\mathrm{K}}: \mathrm{K} \otimes_{\mathbf{Z}} \mathrm{R}_{\mathrm{K}}(\mathrm{A}) \to \mathrm{A}^*\) .

COROLLARY. --- a) Suppose that the field K has characteristic 0. The homomorphisms \(\theta\) and \(\theta_{K}\) are injective. Two semisimple A-modules of finite dimension over K are isomorphic if and only if their characters are equal.

\begin{enumerate}
\def\labelenumi{\alph{enumi})}
\setcounter{enumi}{1}
\tightlist
\item
  Suppose that the field K is perfect of nonzero characteristic p.~Then the homomorphism \(\theta_{K}\) is injective, and the kernel of \(\theta\) is \(p\mathrm{R}_{\mathrm{K}}(\mathrm{A})\) . Let E be a semisimple A-module of finite dimension over K. The linear form \(Tr_{E}\) is zero if and only if there exists an A-module F such that E is isomorphic to \(F^{p}\) .
\end{enumerate}

Let \(S_{K}\) be the set of classes of simple A-modules of finite dimension over K. The elements {[}S{]}, where S runs through \(S_{K}\) , form a basis of the Z-module \(\mathrm{R}_{\mathrm{K}}(\mathrm{A})\) , so the elements \(1 \otimes [S]\) form a basis of the K-vector space \(\mathrm{K} \otimes_{\mathbf{Z}} \mathrm{R}_{\mathrm{K}}(\mathrm{A})\) (VIII, p.~195).

Suppose that the field K is perfect. By Proposition 6 and the corollary of VIII, p.~383, the elements \(\theta([S]) = \theta_{\mathrm{K}}(1 \otimes [S]) = \mathrm{Tr}_{\mathrm{S}}\) of \(A^{*}\) are linearly independent over K. It follows that the homomorphism \(\theta_{K}\) is injective and that the kernel of the homomorphism \(\theta\) consists of the elements \(\sum_{S \in S_{K}} n_{S}[S]\) of \(\mathrm{R}_{\mathrm{K}}(\mathrm{A})\) such that \(n_{S} \cdot 1_{K} = 0\) for every \(S \in S_{K}\) . The kernel of \(\theta\) is therefore equal to \(p R_{\mathrm{K}}(\mathrm{A})\) , where p is the characteristic of K. In particular, if K has characteristic 0, then the homomorphism \(\theta\) is injective.

The last assertion of a) follows from the corollary of VIII, p.~190.

Suppose that the assumptions of b) are satisfied, and let E be a semisimple A-module of finite dimension over K. If there exists an A-module F such that E is the direct sum of p submodules isomorphic to F, then the module F is semisimple and finite-dimensional over K and we have \(Tr_{E} = p Tr_{F} = 0\) . Conversely, suppose \(Tr_{E} = 0\) . We have \([E] \in p R_{K}(A)\) , so the multiplicity of every simple module \(S \in S_{K}\) in E is a multiple of p (VIII, p.~190, Proposition 7); we can write it as \(pn_{S}\) with \(n_{S} \in N\) . The family ( \(n_{S}\) ) has finite support. Set \(F = \oplus_{S \in S_{K}} S^{n_{S}}\) . We then have {[}E{]} = {[}F\^{}\{p\}{]}, so that E and F\^{}\{p\} are isomorphic (VIII, p.~190, Corollary).

Let \(\mathrm{E}\) be an A-module of finite dimension over \(\mathrm{K}\) . Let \(a \in \mathrm{A}\) . Denote by \(\chi_{\mathrm{E}}(a; \mathrm{T})\) the determinant of the endomorphism \(1_{\mathrm{E}} + T a_{\mathrm{E}}\) of the \(\mathrm{K}[\mathrm{T}]\) -module \(\mathrm{E}[\mathrm{T}] = \mathrm{K}[\mathrm{T}] \otimes_{\mathrm{K}} \mathrm{E}\) (III, §8, No.~11, p.~541). We have the relation

\[
\chi_ {\mathrm{E}} (a; \mathrm{T}) \equiv 1 + \mathrm{Tr} _ {\mathrm{E}} (a) \mathrm{T} \pmod {\mathrm{T} ^ {2} \mathrm{K} [ \mathrm{T} ]}\tag{10}
\]

(loc. cit., formula (49)). Since this polynomial has constant term equal to 1, it is an invertible formal power series (IV, §4, No.~4, p.~30). Moreover, if

\[
0 \longrightarrow \mathrm{E} ^ {\prime} \longrightarrow \mathrm{E} \longrightarrow \mathrm{E} ^ {\prime \prime} \longrightarrow 0
\]

is an exact sequence of A-modules, all finite-dimensional over K, then we have \(\chi_{\mathrm{E}}(a;\mathrm{T})=\chi_{\mathrm{E}^{\prime}}(a;\mathrm{T})\chi_{\mathrm{E}^{\prime\prime}}(a,\mathrm{T})\) (III, §8, No.~7, p.~534, formula (31)). By the definition of the Grothendieck group \(\mathrm{R}_{\mathrm{K}}(\mathrm{A})\) (VIII, p.~194) and its universal property (VIII, p.~186, Proposition 4), there exists a unique homomorphism \(\chi_{a}\) from the group \(\mathrm{R}_{\mathrm{K}}(\mathrm{A})\) to the multiplicative group \(1+\mathrm{TK}[[\mathrm{T}]]\) such that \(\chi_{a}([\mathrm{E}])=\chi_{\mathrm{E}}(a;\mathrm{T})\) for every A-module E of finite dimension over K. We deduce from formula (10) the relation

\[
\chi_ {a} (x) \equiv 1 + \theta (x) (a) \mathrm{T} \pmod {\mathrm{T} ^ {2} \mathrm{K} [ [ \mathrm{T} ] ]}\tag{11}
\]

for \(x\in \mathrm{R}_{\mathrm{K}}(\mathrm{A})\) and \(a\in \mathrm{A}\).

If E and \(E'\) are two A-modules of finite dimension over K with isomorphic semisimplifications, then we have \(\chi_{\mathrm{E}}(a;\mathrm{T})=\chi_{\mathrm{E}'}(a;\mathrm{T})\) .

THEOREM 2. --- Let \(\mathcal{A}\) be a generating subset of the K-vector space A. The homomorphism \(\chi_{\mathcal{A}}: \mathrm{R}_{\mathrm{K}}(\mathrm{A}) \to (1 + \mathrm{TK}[[\mathrm{T}]])^{\mathcal{A}}\) defined by the relation \(\chi_{\mathcal{A}}(x) = (\chi_a(x))_{a \in \mathcal{A}}\) is injective.

Let \(x\) be an element of \(\mathrm{R}_{\mathrm{K}}(\mathrm{A})\) such that \(\chi_{\mathcal{A}}(x) = 1\) . By (11), we have \(\theta(x)(a) = 0\) for every \(a \in \mathcal{A}\) , and therefore \(\theta(x) = 0\) because \(\theta(x)\) is a K-linear form on A and \(\mathcal{A}\) generates the K-vector space A. If the characteristic of K is zero, then this implies \(x = 0\) (VIII, p.~384, Corollary of Proposition 6), and therefore the result in this case. Suppose from now on that the characteristic \(p\) of K is nonzero.

Let us first treat the case when K is algebraically closed. By the above and loc. cit., we then have \(x \in p \mathrm{R}_{\mathrm{K}}(\mathrm{A})\) . Let \(y \in \mathrm{R}_{\mathrm{K}}(\mathrm{A})\) be such that x = py. For every element a of A, we have \(\chi_{a}(y)^{p} = \chi_{a}(py) = \chi_{a}(x) = 1\) . So we have \((\chi_{a}(y) - 1)^{p} = 0\) , and therefore \(\chi_{a}(y) = 1\) because the ring K{[}{[}T{]}{]} is an integral domain. So y belongs to the kernel of the endomorphism \(\chi_{A}\) . It follows by induction that x belongs to \(p^{n}\mathrm{R}_{\mathrm{K}}(\mathrm{A})\) for every integer \(n \geqslant 1\) . Since \(\mathrm{R}_{\mathrm{K}}(\mathrm{A})\) is a free Z-module, this implies x = 0, and therefore the injectivity of \(\chi_{A}\) in this case.

If K is no longer supposed to be algebraically closed, then we choose an algebraic closure \(\overline{K}\) of K and consider the diagram of groups and group homomorphisms

\[
\begin{array}{c} \mathrm{R} _ {\mathrm{K}} (\mathrm{A}) \xrightarrow {\chi_ {\mathcal {A}}} (1 + \mathrm{TK} [ [ \mathrm{T} ] ]) ^ {\mathcal {A}} \\ \Biggl \downarrow^ {u} \qquad \qquad \qquad \qquad \Biggl \downarrow^ {i} \\ \mathrm{R} _ {\overline {{\mathrm{K}}}} (\mathrm{A} _ {(\overline {{\mathrm{K}}})}) \xrightarrow {\overline {{\chi}} _ {\mathcal {A}}} (1 + \mathrm{T} \overline {{\mathrm{K}}} [ [ \mathrm{T} ] ]) ^ {\mathcal {A}}, \end{array}\tag{12}
\]

where \(u\) is the homomorphism deduced from the extension of scalars from K to \(\overline{\mathbf{K}}\) (VIII, p.~195), \(i\) is the canonical injection, and \(\overline{\chi}_{\mathcal{A}}\) is the homomorphism \(z \mapsto (\chi_{1 \otimes a}(z))_{a \in \mathcal{A}}\) . By formula (12) of III, §9, No.~1, p.~542, diagram (12) is commutative. By the above, the homomorphism \(\overline{\chi}_{\mathcal{A}}\) is injective. Since \(u\) is injective (VIII, p.~195, Theorem 1), the homomorphism \(\chi_{\mathcal{A}}\) is injective.

COROLLARY 1. --- Let E and F be semisimple A-modules of finite dimension over K, and let \(\mathcal{A}\) be a generating subset of the K-vector space A. Suppose that for every \(a \in \mathcal{A}\) , the characteristic polynomials of the endomorphisms \(a_{E}\) and \(a_{F}\) of the K-vector spaces E and F are equal. Then the A-modules E and F are isomorphic.

Let \(a\) be an element of \(\mathcal{A}\) . The characteristic polynomials of \(a_{\mathrm{E}}\) and \(a_{\mathrm{F}}\) have the same degree, so the dimension of E is equal to that of F; we denote it by \(n\) . Let \(\mathrm{Pc}_{\mathrm{E}}(a;\mathrm{T})\) be the characteristic polynomial of \(a_{\mathrm{E}}\) . In \(\mathrm{K(T)}\) , we have the equalities

\[
\chi_ {\mathrm{E}} (a; \mathrm{T}) = \det \left(1 + a _ {\mathrm{E}} \mathrm{T}\right) = (- \mathrm{T}) ^ {n} \det \left(\frac {- 1}{\mathrm{T}} - a _ {\mathrm{E}}\right) = (- \mathrm{T}) ^ {n} \operatorname{Pc} _ {\mathrm{E}} \left(a; \frac {- 1}{\mathrm{T}}\right),
\]

and \(\chi_{\mathrm{F}}(a;\mathrm{T})\) is given by an analogous formula. Because of our assumptions, we have \(\chi_{\mathrm{E}}(a;\mathrm{T})=\chi_{\mathrm{F}}(a;\mathrm{T})\) . By Theorem 2, we have \([E]=[F]\) , which implies that E and F are isomorphic (VIII, p.~190, Corollary of Proposition 7).

COROLLARY 2. --- Let A be a central simple algebra of finite degree over K. Let B be a semisimple K-algebra, let f and g be algebra homomorphisms from B to A, and let \(\mathcal{B}\) be a generating subset of the K-vector space B. The following properties are equivalent:

\begin{enumerate}
\def\labelenumi{(\roman{enumi})}
\item
  There exists an inner automorphism \(\theta\) of A such that \(g = \theta \circ f\) .
\item
  For every \(b \in \mathcal{B}\) , we have \(\mathrm{Pc}_{\mathrm{A / K}}(f(b); \mathrm{X}) = \mathrm{Pc}_{\mathrm{A / K}}(g(b); \mathrm{X})\) .
\end{enumerate}

When \(\mathrm{K}\) has characteristic zero, these properties are equivalent to the following:

\begin{enumerate}
\def\labelenumi{(\roman{enumi})}
\setcounter{enumi}{2}
\tightlist
\item
  For every \(b \in \mathcal{B}\) , we have \(\mathrm{Tr}_{\mathrm{A / K}}(f(b)) = \mathrm{Tr}_{\mathrm{A / K}}(g(b))\) .
\end{enumerate}

Denote by M and N the left B-modules obtained by endowing the additive group of A with the external laws \((b,a)\mapsto f(b)a\) and \((b,a)\mapsto g(b)a\) , respectively. By Proposition 1 of VIII, p.~257, property (i) is equivalent to the fact that the B-modules M and N are isomorphic. By construction, the homothety \(b_{M}\) associated with an element b of B is left multiplication by \(f(b)\) in A; consequently, we have the relations

\[
\begin{array}{c} \mathrm {Pc_ {M}} (b; \mathrm{X}) = \mathrm {Pc_ {A / K}} (f (b); \mathrm{X}), \\ \mathrm{Tr} _ {\mathrm{M}} (b) = \mathrm{Tr} _ {\mathrm{A/K}} (f (b)) \end{array}
\]

and two analogous relations where we replace M with N and f with g. We know (VIII, p.~386, Corollary 1) that the B-modules M and N are isomorphic if and only if we have \(\mathrm{Pc}_{\mathrm{M}}(b;\mathrm{X})=\mathrm{Pc}_{\mathrm{N}}(b;\mathrm{X})\) for every \(b\in B\) . When the field K has characteristic 0, this relation is also equivalent to \(\mathrm{Tr}_{\mathrm{M}}(b)=\mathrm{Tr}_{\mathrm{N}}(b)\) for every \(b\in B\) (VIII, p.~384, Corollary of Proposition 6). The equivalence of properties (i) and (ii), and also of (i) and (iii) when K has characteristic 0, follows.

